%%%%%%%%%%%%%%%%%%%%%%%%%%%%%%%%%%%%%%%%
%       LISTA DE SÍMBOLOS
%%%%%%%%%%%%%%%%%%%%%%%%%%%%%%%%%%%%%%%%
\chapter*{\hfill{\centering Lista de símbolos}\hfill}
\addcontentsline{toc}{chapter}{Lista de símbolos}
\vspace{0.5cm}
\normalsize

\noindent\textbf{Constantes fundamentales y magnitudes métricas}
\vspace{0.2cm}

\begin{tabular}{ll}
$c$ & --$\quad$ Velocidad de la luz en el vacío ($\text{m}\cdot\text{s}^{-1}$)\\
$G$ & --$\quad$ Constante de gravitación universal ($\text{m}^3\cdot\text{kg}^{-1}\cdot\text{s}^{-2}$)\\
$\varepsilon_0$ & --$\quad$ Permitividad eléctrica del vacío ($\text{F}\cdot\text{m}^{-1}$)\\
$\mu_0$ & --$\quad$ Permeabilidad magnética del vacío ($\text{H}\cdot\text{m}^{-1}$)\\
$M$ & --$\quad$ Masa del agujero negro ($\text{kg}$)\\
$Q$ & --$\quad$ Carga eléctrica neta del agujero negro ($\text{C}$)\\
$J$ & --$\quad$ Momento angular del agujero negro ($\text{kg}\cdot\text{m}^2\cdot\text{s}^{-1}$)\\
$a \equiv \frac{J}{Mc}$ & --$\quad$ Parámetro de espín por unidad de masa ($\text{m}$)\\
$\delta \equiv \frac{Q^2}{4\pi\varepsilon_0 G M^2}$ & --$\quad$ Parámetro adimensional de carga eléctrica débil ($\delta \ll 1$)\\
$r_g \equiv \frac{GM}{c^2}$ & --$\quad$ Radio gravitacional ($\text{m}$)\\
$\theta_g \equiv \frac{r_g}{D}$ & --$\quad$ Escala angular gravitacional ($\text{rad}$ o $\mu\text{as}$)\\
\end{tabular}

\vspace{0.5cm}
\noindent\textbf{Tensores y funciones métricas en $(+,-,-,-)$}
\vspace{0.2cm}

\begin{tabular}{ll}
$g_{\mu\nu}$ & --$\quad$ Tensor métrico del espaciotiempo con firma $(+,-,-,-)$\\
$g \equiv \det(g_{\mu\nu})$ & --$\quad$ Determinante del tensor métrico ($g < 0$)\\
$R_{\mu\nu}$ & --$\quad$ Tensor de curvatura de Ricci\\
$R$ & --$\quad$ Escalar de curvatura de Ricci ($R = g^{\mu\nu} R_{\mu\nu}$)\\
$F_{\alpha\beta}$ & --$\quad$ Tensor de fuerza del campo electromagnético\\
$A_\mu$ & --$\quad$ Cuadripotencial electromagnético ($\text{V}\cdot\text{s}\cdot\text{m}^{-1}$)\\
$T_{\mu\nu}^{(\mathrm{EM})}$ & --$\quad$ Tensor de energía-impulso electromagnético ($\text{J}\cdot\text{m}^{-3}$)\\
$\Sigma(r,\theta)$ & --$\quad$ Función métrica auxiliar de Boyer-Lindquist: $r^2 + a^2 \cos^2\theta$ ($\text{m}^2$)\\
$\Delta(r)$ & --$\quad$ Función de horizonte de Kerr-Newman ($\text{m}^2$)\\
\end{tabular}
\newpage
\vspace{0.5cm}
\noindent\textbf{Dinámica de fotones, constantes de movimiento y teoría perturbativa}
\vspace{0.2cm}

\begin{tabular}{ll}
$S$ & --$\quad$ Función de acción principal de Hamilton-Jacobi ($\text{J}\cdot\text{s}$)\\
$E$ & --$\quad$ Energía del fotón conservada por simetría temporal ($\text{J}$)\\
$L_z$ & --$\quad$ Momento angular azimutal conservado por simetría axial ($\text{kg}\cdot\text{m}^2\cdot\text{s}^{-1}$)\\
$\mathcal{K}$ & --$\quad$ Constante de separación de Carter ($\text{J}^2\cdot\text{m}^2\cdot\text{s}^{-2}$)\\
$\mathcal{R}(r)$ & --$\quad$ Potencial radial efectivo de Hamilton-Jacobi ($\text{J}^2\cdot\text{m}^2$)\\
$\Phi \equiv \frac{c L_z}{E}$ & --$\quad$ Parámetro de impacto azimutal reducido ($\text{m}$)\\
$\mathcal{Q} \equiv \frac{c^2 \mathcal{K}}{E^2}$ & --$\quad$ Parámetro de impacto polar reducido de Carter ($\text{m}^2$)\\
$b_c \equiv \sqrt{\mathcal{Q}_c + \Phi_c^2}$ & --$\quad$ Parámetro de impacto crítico ($\text{m}$)\\
$r_{\mathrm{ph}}$ & --$\quad$ Radio de la órbita fotónica esférica inestable ($\text{m}$)\\
$R_{\mathrm{eff}}(\delta)$ & --$\quad$ Radio efectivo equivalente de la sombra ($\text{m}$ o $\mu\text{as}$)\\
$D = 1/27$ & --$\quad$ Coeficiente perturbativo de segundo orden del área de la sombra\\
$\gamma_\phi$ & --$\quad$ Exponente orbital de Lyapunov\\
$f_{\mathrm{thin}} \equiv e^{-2\pi\gamma_\phi}$ & --$\quad$ Factor de adelgazamiento exponencial de Sneppen\\
\end{tabular}

\vspace{0.5cm}
\noindent\textbf{Interferometría de muy larga base (VLBI) y observables EHT}
\vspace{0.2cm}

\begin{tabular}{ll}
$(\alpha_c, \beta_c)$ & --$\quad$ Coordenadas celestes en el plano del observador ($\mu\text{as}$)\\
$(u, v)$ & --$\quad$ Coordenadas de frecuencia espacial en el plano de Fourier ($\text{G}\lambda$)\\
$u \equiv \sqrt{u^2 + v^2}$ & --$\quad$ Distancia escalar de la línea de base ($\text{G}\lambda$)\\
$V(u, v)$ & --$\quad$ Visibilidad compleja interferométrica ($\text{Jy}$)\\
$I(\alpha_c, \beta_c)$ & --$\quad$ Distribución de brillo específico en el cielo ($\text{W}\cdot\text{m}^{-2}\cdot\text{Hz}^{-1}\cdot\text{sr}^{-1}$)\\
$F_0$ & --$\quad$ Flujo de emisión total no polarizado ($\text{Jy}$)\\
$w$ & --$\quad$ Ancho FWHM del perfil de emisión del anillo ($\mu\text{as}$)\\
$J_0(x)$ & --$\quad$ Función de Bessel de primer tipo y orden cero\\
$u_1$ & --$\quad$ Frecuencia espacial del primer nulo de visibilidad ($\text{G}\lambda$)\\
$\chi^2_\nu$ & --$\quad$ Chi-cuadrado reducido para la evaluación de bondad de ajuste\\
$\Delta\chi^2_\nu$ & --$\quad$ Diferencia de bondad de ajuste respecto al modelo de Schwarzschild\\
\end{tabular}